\documentclass[]{article}  

% Default font is the helvetica postscript font
\usepackage{helvet}
\usepackage{hyperref}
\usepackage{amsmath}
\usepackage{amsthm}
\usepackage{amssymb}
\usepackage[a4paper,margin=25mm]{geometry}
\usepackage{fancyhdr}
\usepackage{lipsum}
\usepackage[ruled,vlined]{algorithm2e}
\usepackage[utf8]{inputenc}
\usepackage[english]{babel}
\usepackage{graphicx}


\newtheorem{theorem}{Theorem}
\newtheorem{lemma}[theorem]{Lemma}
\newtheorem{defn}[theorem]{Definition}
\newtheorem{remark}[theorem]{Remark}
%\renewcommand{\baselinestretch}{1.15}
\newcommand{\s}{\mathbf{s}}
\newcommand{\ba}{\mathbf{a}}
\newcommand{\bw}{\mathbf{w}}
\newcommand{\bx}{\mathbf{x}}
\newcommand{\bv}{\mathbf{v}}
\newcommand{\thetab}{\boldsymbol{{\theta}}}
\newcommand{\h}{\mathbf{h}}

\newcommand{\cA}{\mathcal{A}}
\newcommand{\cW}{\mathcal{W}}
\newcommand{\cH}{\mathcal{H}}

\newcommand{\R}{\mathbb{R}}
\newcommand{\E}{\mathbb{E}}
\newcommand{\Prob}{\mathbb{P}}
\newcommand{\KL}{\mathrm{D}_{KL}}



\title{Learnability of a Finite Hypothesis Class in Learning Problems with Noisy Training-Set Labels}
\author{Behrad Moniri\\\texttt{bemoniri@live.com}}
\date{}

\usepackage{sectsty} \sectionfont{\fontsize{15}{15}\selectfont}
\begin{document}
\setcounter{section}{-1}
\maketitle

\begin{abstract}
In this note, the PAC-Learnability of a finite hypothesis class is proved for a learning problem in which the training set labels are flipped with a certain probability.
\end{abstract}

%\tableofcontents
Let $\cH$ be a finite hypothesis class, $h^*\in \cH$ be the target concept and $S = \{\bx_1, \bx_2, \dots, \bx_m\}$ be the training set. The training set labels are corrupted with noise, i.e., the label $y_i$ is reported as $h^*(\bx_i) \oplus \zeta$ in which $\zeta \sim \text{Bern }(\eta)$.  We will prove that the algorithm which chooses the function $h_S$ with the least number of errors on training points, has a low generalization error with high probability.

\begin{defn}
For each $h\in \cH$ and sample set $S$, Define $d(h)$ as the probability that the label $y$ of a training sample $\bx$ is not not equal to $h^*(\bx)$. Also define $\hat{d}(h)$ as the empirical loss of $h$ on the training set, 
   $\hat{d}(h) = \frac{\sum_{i = 1}^{m} \mathbf{1}[y_i \neq h(\bx_i) ]}{m}$.
\end{defn}

\begin{remark}
Trivially we have $d(h^*) = \eta$.
\end{remark}


We will now prove a very useful lemma linking $d(h)$ and the generalization error of $h$.
\begin{lemma}
\label{lem}
Let $\epsilon>0$ be an arbitrary given number. If $L_D(h)>\epsilon$, then we have $d(h) - d(h^*) \geq \epsilon'$ in which $\epsilon' = \epsilon(1-2\eta)$.
\end{lemma}

\begin{proof} 
The label of a training set is flipped with probability $\eta$. Hence, we have
\begin{align}
d(h) &= \eta \;\mathbb{P}[h^*(\bx) = h(\bx)] + (1-\eta)\; \mathbb{P}[h^*(\bx) \neq h(\bx)]\\
&= \eta\big(1-L_D(h)\big) + \big(1-\eta\big) L_D(h)\\
&= \eta + (1-2\eta)L_D(h).
\end{align} 
For a given $h$, if $L_D(h)>\epsilon$, we have $d(h) - \eta = d(h) - d(h^*) > (1-2\eta) \epsilon = \epsilon'$.
\end{proof}
 We now restate the well known Hoeffding's inequality (Theorem \eqref{Hoeffding}) as it is used in the proof of our main theorem.
\begin{theorem}
\label{Hoeffding}
Let $X_1, X_2, \dots, X_m$ be independent random variables. Also assume that each $X_i$ takes values in $[a_i, b_i]$ with probability 1. For any $\epsilon>0$, the following inequalities hold for $S_m = \sum_{i = 1}^{m}X_i$:
\begin{align}
\mathbb{P}\big[S_m - \mathbb{E}[S_m] \geq \epsilon\big] &\leq e^{-2\epsilon^2/\sum_{i = 1}^{m}(b_i-a_i)^2},\\
\mathbb{P}\big[S_m - \mathbb{E}[S_m] \leq -\epsilon\big] &\leq e^{-2\epsilon^2/\sum_{i = 1}^{m}(b_i-a_i)^2}.
\end{align}
\end{theorem} 
\begin{proof}
See Appendix (D) of \cite{mohri}.
\end{proof}
We are now ready to state and prove our main claim.
\begin{theorem}
In the noisy label setting with noise parameters $\eta$ and $\eta$ with any distribution $D$, let $\epsilon, \delta >0$ be arbitrary given numbers. Given 
$m\geq \frac{2}{\epsilon^2(1-2\eta)^2}\Big(\log(|\cH|) +  \log(\frac{2}{\delta})\Big)$ training samples, for any $h \in \cH$  with $L_D(h) >\epsilon$, we have 
\begin{equation}
\hat{d}(h) - \hat{d}(h^*)  \geq 0
\end{equation}
with probability at least $\delta$.
\end{theorem}
\begin{proof}
Let $h\in\cH$ be any hypothesis with $L_D(h) > \epsilon$. We can decompose $\hat{d}(h) - \hat{d}(h^*)$ as follows:
\begin{align}
\label{decomp}
\hat{d}(h) - \hat{d}(h^*) = [\hat{d}(h) - d(h)]+[d(h) - d(h^*)]+[d(h^*) - \hat{d}(h^*)].
\end{align}
Based on lemma \eqref{lem}, we know that for any $h$ with $L_D(h)>\epsilon$, we have $d(h)-d(h^*)> \epsilon'$.  We will now show that given enough samples, each of the remaining terms are greater than $\frac{-\epsilon'}{2}$ with high probability.

\begin{enumerate}
\item First, we will show that given $m\geq \frac{2}{\epsilon^2(1-2\eta)^2}\Big(\log(\frac{2}{\delta})\Big)$ samples, we have $\hat{d}(h^*) - d(h^*)  \leq \frac{\epsilon'}{2}$ with probability at least $1-\delta/2$. This is a direct consequence of the Hoeffding's inequality. Note that $\mathrm{E}[\hat{d}(h^*)] = d(h^*)$ and we have 
\begin{equation}
\hat{d}(h) = \frac{\sum_{i = 1}^{m} \mathbf{1}[y_i \neq h(\bx_i) ]}{m} = \sum_{i = 1}^{m} X_i,
\end{equation}
 in which $X_i \in [0, \frac{1}{m}]$ with probability 1 and $X_i$s are jointly independent. The Hoeffding's inequality yields
\begin{align}
\mathbb{P}\big[\hat{d}(h^*) - d(h^*) \geq \frac{\epsilon'}{2}\big] &\leq \exp\Big[{-\frac{m\epsilon'^2}{2}}\Big] \leq \frac{\delta}{2}.
\end{align}
Where the last inequality follows from $m\geq \frac{2}{\epsilon^2(1-2\eta)^2}\Big(\log(\frac{2}{\delta})\Big)$.
\item Second, we will prove a \textit{Uniform Convergence} property for $\cH$. We will prove that given $m\geq \frac{2}{\epsilon^2(1-2\eta)^2}\Big(\log(|\cH|) +  \log(\frac{2}{\delta})\Big)$ samples, the event
\begin{equation}
\forall h\in \cH,\; d(h) - \hat{d}(h) \leq \frac{\epsilon'}{2}
\end{equation}
occurs with probability at least  $1-\frac{\delta}{2}$. 

To prove it, we can write the following chain of inequalities:
\begin{align}
\mathbb{P}\Big[\forall h\in \cH, d(h) - \hat{d}(h) \geq \frac{\epsilon'}{2}\Big] &= \mathbb{P}\Big[\bigcup_{h\in\cH} \big\{d(h) - \hat{d}(h) \geq \frac{\epsilon'}{2}\big\}\Big]\\
&\leq \sum_{h \in \cH}\mathbb{P}\Big[d(h) - \hat{d}(h) \geq \frac{\epsilon'}{2}\Big]\\
&\leq \sum_{h \in \cH} \exp\Big[{-\frac{m\epsilon'^2}{2}}\Big] = |\cH|\exp\Big[{-\frac{m\epsilon'^2}{2}}\Big]
\end{align}
Given $m\geq \frac{2}{\epsilon^2(1-2\eta)^2}\Big(\log(\frac{2}{\delta})+\log(|\cH|)\Big)$, the right hand side is less than or equal to $\frac{\delta}{2}$.
\end{enumerate}

\noindent
Thus, given $m\geq \frac{2}{\epsilon^2(1-2\eta)^2}\Big(\log(\frac{2}{\delta})+\log(|\cH|)\Big)$, the first term and third terms of \eqref{decomp} are both greater than $-\frac{\epsilon'}{2}$ with probability at least $1-\frac{\delta}{2}$.  Hence, with probability at least $1-\delta$, 
\begin{align}
\hat{d}(h) - \hat{d}(h^*) &= [\hat{d}(h) - d(h)]+[d(h) - d(h^*)]+[d(h^*) - \hat{d}(h^*)]\\
&> -\frac{\epsilon'}{2} + \epsilon' - \frac{\epsilon'}{2} = 0,
\end{align}
which proves the theorem.
\end{proof}


This theorem states that given $m\geq \frac{2}{\epsilon^2(1-2\eta)^2}\Big(\log(\frac{2}{\delta})+\log(|\cH|)\Big)$ samples, the probability of the algorithm $h_S = \arg\min_{h\in\cH} \hat{d}(h)$ having $L_D(h_S)\leq\epsilon$ is at least than $1-\delta$.


\bibliography{bib}
\bibliographystyle{alpha}


\end{document}